\documentclass[11pt]{article}

\usepackage[margin=1in]{geometry}
\usepackage[T1]{fontenc}
\usepackage{lmodern}
\usepackage{microtype}
\usepackage{amsmath,amssymb,amsthm,mathtools}
\usepackage{enumitem}
\usepackage[hidelinks]{hyperref}

\newtheorem{theorem}{Theorem}[section]
\newtheorem{proposition}[theorem]{Proposition}
\newtheorem{lemma}[theorem]{Lemma}
\newtheorem{corollary}[theorem]{Corollary}
\theoremstyle{definition}
\newtheorem{definition}[theorem]{Definition}
\theoremstyle{remark}
\newtheorem{remark}[theorem]{Remark}

\newcommand{\N}{\mathbb N}
\newcommand{\K}{\mathbb K}
\newcommand{\F}{\mathcal F}
\newcommand{\G}{\mathcal G}
\newcommand{\Ccal}{\mathcal C}
\newcommand{\U}{\mathbb U}
\newcommand{\Aut}{\operatorname{Aut}}
\newcommand{\Isom}{\operatorname{Isom}}
\newcommand{\Sym}{\operatorname{Sym}}
\newcommand{\supp}{\operatorname{supp}}
\newcommand{\preF}{\preccurlyeq_{\F}}
\newcommand{\simF}{\sim_{\F}}

\title{Replacement Classes and the Isometry Groups\\
of Combinatorial Banach Spaces}
\author{A research note on the questions of Brech--Ferenczi--Tcaciuc}
\date{21 June 2026}

\begin{document}
\maketitle

\begin{abstract}
For a family $\F\subseteq[\N]^{<\omega}$, define the replacement relation
$i\preccurlyeq_{\F}j$ by requiring that $j$ may replace $i$ in every finite
context without destroying membership in $\F$.  We prove that this relation
is a preorder for every family and is total whenever $\F$ is spreading.  Its
equivalence classes are consecutive intervals of $\N$, and the full
permutation automorphism group of a spreading family is exactly the
unrestricted product of the symmetric groups on those intervals.  Combining
this with the signed-permutation theorem of Brech--Ferenczi--Tcaciuc gives a
complete intrinsic description of the surjective linear isometry group of
every regular combinatorial Banach space and of its dual.  We also classify
which permutation groups occur, verify the supplied threshold-family result,
and obtain a short rigidity theorem: two permutation-isomorphic spreading
families on $\N$ are equal.  A literature search found close finite
antecedents through desirability and vicinal preorders, but no exact statement
of the countably infinite Banach-space conclusion; this bibliographic
observation is not a formal novelty certification.
\end{abstract}

\section{Problem, verdict on the partial result, and main conclusion}

Throughout, $\N=\{1,2,\ldots\}$.  A family means a subset of
$[\N]^{<\omega}$.  It is \emph{spreading} if every coordinatewise increase of
a member remains in the family.  It is \emph{regular} if it is hereditary,
spreading, compact as a subspace of $2^{\N}$, and contains all singletons.
For a regular family $\F$, the associated combinatorial Banach space is the
completion of $c_{00}$ under
\[
  \|x\|_{\F}=\sup_{A\in\F}\sum_{n\in A}|x(n)|.
\]

Brech, Ferenczi, and Tcaciuc proved that every surjective linear isometry of a
combinatorial space, or of its dual, is a signed permutation of the canonical
basis.  They then asked which permutations preserve a given regular family
and for which combinatorial spaces the full isometry group can be explicitly
described \cite{BFT}.  The same problem was still listed in the 2022 survey of
Antunes and Beanland \cite[Problem~3.1]{ABsurvey}.

The supplied partial result treats
\[
  \F_g=\{\varnothing\}\cup
  \bigl\{A\in[\N]^{<\omega}:A\ne\varnothing,
      \ |A|\le g(\min A)\bigr\},
\]
where $g:\N\to\N$ is nondecreasing and $g\ge1$.  Its conclusion is that
$\Aut(\F_g)$ consists precisely of the permutations preserving every level
set of $g$.

\begin{remark}[Verdict on the supplied argument]
The partial result is correct.  In particular, its compactness proof, the
invariant
\[
  h_{\F_g}(n)=\max\{|A|:A\in\F_g,\ n\in A\}=g(n),
\]
and the converse argument for level-preserving permutations are all sound.
Section~\ref{sec:threshold} gives a self-contained verification.  The main
point of the present note is that the same replacement idea, formulated at
the right level of generality, settles the permutation problem for every
spreading family.
\end{remark}

Here is the resulting answer.  The canonical partition appearing in the
statement is defined in Section~\ref{sec:preorder}.

\begin{theorem}[Main isometry theorem]\label{thm:main-isometry}
Let $\F$ be a regular family, let $\Ccal_{\F}$ be its replacement classes, and
let $\U_{\K}=\{-1,1\}$ in the real case and
$\U_{\K}=\{z\in\mathbb C:|z|=1\}$ in the complex case.  A map
$T:X_{\F}\to X_{\F}$ is a surjective linear isometry if and only if there are
$\varepsilon=(\varepsilon_n)_{n\in\N}\in\U_{\K}^{\N}$ and a permutation
$\pi$ satisfying $\pi(C)=C$ for every $C\in\Ccal_{\F}$ such that
\[
  T e_n=\varepsilon_n e_{\pi(n)}\qquad(n\in\N).
\]
The analogous statement holds on $X_{\F}^{*}$.  Consequently,
\[
  \Isom(X_{\F})\cong
  \U_{\K}^{\N}\rtimes
  \prod_{C\in\Ccal_{\F}}\Sym(C),
\]
with the same permutation part for $\Isom(X_{\F}^{*})$.
\end{theorem}

The product in this formula is the unrestricted direct product: independent
permutations may be chosen on all replacement classes, including an infinite
class.  The classes are consecutive intervals, so the permutation part is a
Young subgroup associated with an interval partition of $\N$.

\section{The replacement preorder}\label{sec:preorder}

The relation below is the set-system analogue of Isbell's desirability
relation for simple games \cite{Isbell}; closely related generalized vicinal
preorders occur in the theory of shifted simplicial complexes
\cite{Klivans}.  We include all details because two elementary features of
the relation do the entire job.

\begin{definition}[Replacement preorder]\label{def:replacement}
For a family $\F\subseteq[\N]^{<\omega}$ and $i,j\in\N$, write
$i\preccurlyeq_{\F}j$ if
\[
  A\cup\{i\}\in\F\quad\Longrightarrow\quad A\cup\{j\}\in\F
  \tag{2.1}\label{eq:replacement}
\]
for every finite $A\subseteq\N\setminus\{i,j\}$.  Thus $j$ can replace $i$
in every context.  Write
\[
  i\sim_{\F}j
  \quad\Longleftrightarrow\quad
  i\preccurlyeq_{\F}j\text{ and }j\preccurlyeq_{\F}i.
\]
The equivalence classes of $\sim_{\F}$ will be called the
\emph{replacement classes} of $\F$ and denoted by $\Ccal_{\F}$.
\end{definition}

\begin{proposition}\label{prop:preorder}
For every family $\F$, the relation $\preccurlyeq_{\F}$ is a preorder.  If
$\pi$ is a permutation and $\G=\pi[\F]$, then
\[
  i\preccurlyeq_{\F}j
  \quad\Longleftrightarrow\quad
  \pi(i)\preccurlyeq_{\G}\pi(j).
  \tag{2.2}\label{eq:transport}
\]
\end{proposition}

\begin{proof}
Reflexivity is immediate.  For transitivity, the cases with repeated indices
are immediate, so suppose that $i,j,k$ are pairwise distinct,
$i\preccurlyeq_{\F}j$, and $j\preccurlyeq_{\F}k$.  Let
$A\subseteq\N\setminus\{i,k\}$ be finite and suppose
$A\cup\{i\}\in\F$.

If $j\notin A$, replace $i$ by $j$ and then $j$ by $k$.  This gives
$A\cup\{k\}\in\F$.  If $j\in A$, put $B=A\setminus\{j\}$.  Starting from
$B\cup\{i,j\}\in\F$, first use $j\preccurlyeq_{\F}k$ with context
$B\cup\{i\}$ to obtain $B\cup\{i,k\}\in\F$.  Then use
$i\preccurlyeq_{\F}j$ with context $B\cup\{k\}$ to obtain
$B\cup\{j,k\}=A\cup\{k\}\in\F$.  Thus
$i\preccurlyeq_{\F}k$.

For \eqref{eq:transport}, apply $\pi^{-1}$ to a context avoiding
$\pi(i),\pi(j)$, use \eqref{eq:replacement} in $\F$, and apply $\pi$ again.
The reverse implication follows in the same way from $\pi^{-1}$.
\end{proof}

\begin{proposition}\label{prop:spreading-structure}
If $\F$ is spreading, then:
\begin{enumerate}[label=\textup{(\roman*)},leftmargin=2.2em]
  \item $i<j$ implies $i\preccurlyeq_{\F}j$, so the replacement preorder is
        total;
  \item each replacement class is a consecutive interval of $\N$;
  \item for $i\ne j$, one has $i\sim_{\F}j$ if and only if the transposition
        $(i\ j)$ belongs to $\Aut(\F)$.
\end{enumerate}
\end{proposition}

\begin{proof}
For (i), if $A\cup\{i\}\in\F$ and $A$ avoids $i,j$, the map that fixes $A$
and sends $i$ to $j$ is an injective pointwise increase.  The spreading
property therefore gives $A\cup\{j\}\in\F$.

For (ii), suppose $i<k<j$ and $i\sim_{\F}j$.  By (i),
$i\preccurlyeq_{\F}k\preccurlyeq_{\F}j$.  Combining these inequalities with
$j\preccurlyeq_{\F}i$ and using transitivity gives both
$i\sim_{\F}k$ and $k\sim_{\F}j$.

For (iii), if the transposition is an automorphism, applying it to
$A\cup\{i\}$ or $A\cup\{j\}$ proves the two replacement implications.
Conversely, if $i\sim_{\F}j$, the transposition fixes every finite set that
contains both or neither of $i,j$; when exactly one is present, the relevant
replacement implication preserves membership.  Hence the transposition
preserves $\F$ in both directions.
\end{proof}

The quotient $\N/{\sim_{\F}}$ is ordered by
\[
  [i]_{\F}\le [j]_{\F}
  \quad\Longleftrightarrow\quad i\preccurlyeq_{\F}j.
  \tag{2.3}\label{eq:quotient-order}
\]
This is well defined by Proposition~\ref{prop:preorder} and is a linear order
when $\F$ is spreading.  More is true: it is a well-order.  Indeed, for any
nonempty collection of classes, the class containing the least integer in
their union is its least member.  Equivalently, the replacement classes form
either a finite interval partition of $\N$ or a sequence of finite consecutive
intervals indexed by $\N$.

\section{Automorphisms of every spreading family}

\begin{theorem}[Automorphism theorem]\label{thm:aut}
Let $\F\subseteq[\N]^{<\omega}$ be spreading.  Then
\[
  \Aut(\F)
  =\bigl\{\pi\in\Sym(\N):\pi(C)=C
       \text{ for every }C\in\Ccal_{\F}\bigr\}
  \cong\prod_{C\in\Ccal_{\F}}\Sym(C).
  \tag{3.1}\label{eq:aut-product}
\]
In particular, the replacement classes are exactly the orbits of
$\Aut(\F)$ on $\N$.
\end{theorem}

\begin{proof}
Let $\pi\in\Aut(\F)$.  Proposition~\ref{prop:preorder} shows that $\pi$
preserves the replacement preorder and hence induces an order automorphism
of the well-ordered quotient $\N/{\sim_{\F}}$.  A well-order has no
nonidentity order automorphism, so every replacement class is fixed setwise.
This proves the inclusion from left to right in \eqref{eq:aut-product}.

Conversely, suppose that $\pi(C)=C$ for every replacement class.  Fix a finite
set $A\subseteq\N$.  For each of the finitely many classes $C$ meeting $A$,
the finite partial bijection
\[
  \pi\big|_{A\cap C}:A\cap C\longrightarrow \pi(A\cap C)
\]
extends to a permutation $\sigma_C$ of $C$ with finite support.  One explicit
construction is to view the partial bijection on the finite set
$(A\cap C)\cup\pi(A\cap C)$ as a disjoint union of directed paths and cycles
and close each path into a cycle.  Let $\sigma$ agree with these $\sigma_C$
and be the identity elsewhere.  Then $\sigma$ has finite support and is a
finite product of transpositions inside replacement classes.  By
Proposition~\ref{prop:spreading-structure}(iii), every such transposition is
an automorphism of $\F$, so $\sigma\in\Aut(\F)$.  Since
$\sigma(A)=\pi(A)$,
\[
  A\in\F\quad\Longleftrightarrow\quad\pi(A)\in\F.
\]
This holds for every finite $A$, hence $\pi\in\Aut(\F)$.

Finally, arbitrary permutations inside a class are allowed, so each class is
one orbit, while the first part shows that no orbit meets two classes.
\end{proof}

\begin{corollary}[Block-count normal form]\label{cor:block-count}
Let $\F$ be spreading.  If finite sets $A,B\subseteq\N$ satisfy
\[
  |A\cap C|=|B\cap C|\qquad(C\in\Ccal_{\F}),
\]
then $A\in\F$ if and only if $B\in\F$.
\end{corollary}

\begin{proof}
Choose, independently in each replacement class, a permutation sending
$A\cap C$ to $B\cap C$, and combine them.  The resulting block-preserving
permutation is an automorphism by Theorem~\ref{thm:aut}.
\end{proof}

Thus every spreading family is completely symmetric inside its replacement
blocks.  All asymmetry is encoded by the ordered list of blocks and by which
finite block-count vectors are admitted.

\section{Consequences for combinatorial Banach spaces}

We now specialize to regular families and use the signed-permutation theorem
of Brech--Ferenczi--Tcaciuc.  Their result says that every surjective linear
isometry between dual combinatorial spaces has the form
$e_n^*\mapsto\varepsilon_n e_{\pi(n)}^*$, and their corollary says that the
permutation $\pi$ is admissible exactly when it carries one defining family
onto the other.  The corresponding assertion for the primal spaces follows
by adjoints \cite[Theorem~10 and Corollary~12]{BFT}.

\begin{proof}[Proof of Theorem~\ref{thm:main-isometry}]
By the cited theorem, a surjective linear isometry of $X_{\F}$ or
$X_{\F}^{*}$ is a signed permutation, and its permutation part belongs to
$\Aut(\F)$.  Theorem~\ref{thm:aut} identifies this automorphism group with
$\prod_{C\in\Ccal_{\F}}\Sym(C)$.  Conversely, every coordinatewise sign
change is an isometry because the canonical basis is $1$-unconditional, and
every block-preserving permutation preserves $\F$, hence the defining norm.
The displayed semidirect product uses the natural action of coordinate
permutations on sign sequences.
\end{proof}

This gives a universal intrinsic answer to the two questions in
\cite{BFT}: for every regular family, compute the replacement equivalence
classes, and then permit arbitrary permutations within each class and no
permutations between classes.

\section{Verification and extension of the threshold-family result}
\label{sec:threshold}

\begin{proposition}[Threshold families]\label{prop:threshold}
Let $g:\N\to\N$ be nondecreasing with $g(n)\ge1$, and define
\[
  \F_g=\{\varnothing\}\cup
  \{A\in[\N]^{<\omega}:A\ne\varnothing,
       \ |A|\le g(\min A)\}.
\]
Then $\F_g$ is regular and
\[
  \Aut(\F_g)
  =\{\pi\in\Sym(\N):g(\pi(n))=g(n)\text{ for all }n\}.
  \tag{5.1}\label{eq:threshold-aut}
\]
Its replacement classes are exactly the nonempty level sets
$L_m=\{n:g(n)=m\}$.
\end{proposition}

\begin{proof}
If $B\subseteq A\in\F_g$ is nonempty, then
$|B|\le |A|\le g(\min A)\le g(\min B)$, so $\F_g$ is hereditary.  If
$A=\{n_1<\cdots<n_k\}\in\F_g$ and $m_r\ge n_r$ with
$m_1<\cdots<m_k$, then $k\le g(n_1)\le g(m_1)$, so $\F_g$ is spreading.
All singletons belong to it.

For compactness, it is enough to prove that $\F_g$ is closed in the compact
metrizable space $2^{\N}$.  Suppose $A_r\in\F_g$ converges pointwise to
$A\subseteq\N$.  If $A$ were infinite and $a=\min A$, then eventually $A_r$
would contain the first $g(a)+1$ members of $A$ (including $a$), while
containing no integer below $a$.  Hence $\min A_r=a$ and $|A_r|\ge g(a)+1$, a
contradiction.  Thus $A$ is finite.  If $A\ne\varnothing$ and $a=\min A$,
then eventually every member of $A$ lies in $A_r$ and no integer below $a$
does.  Consequently
$|A|\le |A_r|\le g(a)$, so $A\in\F_g$.  The empty set is already included.

For $n\in\N$, set
\[
  h_{\F_g}(n)=\max\{|A|:A\in\F_g,\ n\in A\}.
\]
Every member containing $n$ has size at most
$g(\min A)\le g(n)$, while
$\{n,n+1,\ldots,n+g(n)-1\}\in\F_g$.  Therefore
$h_{\F_g}(n)=g(n)$.  Since a family automorphism preserves the cardinalities
of members containing a given point, every automorphism preserves $g$.

Conversely, suppose $g\circ\pi=g$.  If $A\in\F_g$ is nonempty, let
$r=\min A$, and choose $s\in A$ with $\pi(s)=\min\pi(A)$.  Since $s\ge r$,
\[
  |\pi(A)|=|A|\le g(r)\le g(s)=g(\pi(s))=g(\min\pi(A)).
\]
Thus $\pi(A)\in\F_g$, and the same argument for $\pi^{-1}$ gives equality of
the families.  This proves \eqref{eq:threshold-aut}.

The level sets are the orbits of this automorphism group.  By
Theorem~\ref{thm:aut}, the orbits are exactly the replacement classes.
\end{proof}

The proposition strictly contains the step-function example in
\cite[Example~18]{BFT}: the threshold may have arbitrary finite or infinite
plateaus and may jump by more than one.  Constant $g$ gives the full symmetric
permutation group, while $g(n)=n$ gives the first Schreier family and only
singleton replacement classes.

\section{Exactly which permutation groups occur?}

Call $(I_r)_{r\in R}$ an \emph{interval partition} of $\N$ if $R$ is either
$\{1,\ldots,q\}$ or $\N$, the sets $I_r$ are nonempty consecutive intervals,
$I_r<I_s$ for $r<s$, and $\bigcup_r I_r=\N$.

\begin{theorem}[Realization and classification]\label{thm:realization}
As concrete subgroups of $\Sym(\N)$, the permutation parts of the isometry
groups of regular combinatorial spaces are exactly the interval Young
subgroups
\[
  \prod_{r\in R}\Sym(I_r)
  \tag{6.1}\label{eq:interval-young}
\]
associated with interval partitions of $\N$.
\end{theorem}

\begin{proof}
For every regular family, Proposition~\ref{prop:spreading-structure} and
Theorem~\ref{thm:aut} show that its automorphism group has the form
\eqref{eq:interval-young}.

Conversely, fix an interval partition $(I_r)_{r\in R}$ and define
$g(n)=r$ for $n\in I_r$.  Then $g$ is nondecreasing and at least one.
Proposition~\ref{prop:threshold} shows that $\F_g$ is regular, its replacement
classes are precisely the $I_r$, and its automorphism group is exactly
\eqref{eq:interval-young}.
\end{proof}

If the partition has finitely many blocks, its last block is necessarily an
infinite tail.  If it has infinitely many blocks, every block is finite.  The
full isometry group is obtained from \eqref{eq:interval-young} by adjoining
all coordinatewise signs as in Theorem~\ref{thm:main-isometry}.

\section{Permutation rigidity of spreading families}

Brech and Pi\~na proved that hereditary spreading families on $\N$ which are
images of one another under a permutation must be equal
\cite[Theorem~14]{BP}.  The replacement-class theorem gives a short proof and
shows that heredity is unnecessary.

\begin{theorem}[Rigidity]\label{thm:rigidity}
Let $\F,\G\subseteq[\N]^{<\omega}$ be spreading.  If
$\G=\pi[\F]$ for a permutation $\pi$ of $\N$, then $\F=\G$.
\end{theorem}

\begin{proof}
By Proposition~\ref{prop:preorder}, $\pi$ transports the replacement preorder
of $\F$ to that of $\G$.  Hence it maps the first replacement class of $\F$
onto the first replacement class of $\G$, the second onto the second, and so
on, preserving the cardinality of each class.

An ordered interval partition of $\N$ is determined by this ordered list of
cardinalities.  Indeed, the first blocks are initial intervals of the same
size and hence coincide.  Inductively, after the preceding blocks have been
identified, the next blocks begin at the same least unused integer and have
the same size, so they coincide as well.  If one block is infinite, it is the
final tail and the argument terminates.

Thus $\F$ and $\G$ have the same replacement classes, and $\pi$ preserves
each of them setwise.  Theorem~\ref{thm:aut} then gives
$\pi\in\Aut(\F)$.  Therefore $\G=\pi[\F]=\F$.
\end{proof}

\begin{corollary}
If $\F$ and $\G$ are regular and $X_{\F}$ is linearly isometric to
$X_{\G}$ (equivalently, their duals are linearly isometric), then
$\F=\G$.  After this identification, every isometry has the blockwise signed
form of Theorem~\ref{thm:main-isometry}.
\end{corollary}

\begin{proof}
Brech--Ferenczi--Tcaciuc reduce an isometry to a permutation satisfying
$\G=\pi[\F]$ \cite[Corollary~12]{BFT}; apply
Theorem~\ref{thm:rigidity}.  This also recovers the regular-family
Banach--Stone conclusion of \cite[Corollary~15]{BP}.
\end{proof}

\section{Literature status and scope}

The replacement relation itself is not new in spirit.  Isbell introduced the
desirability relation for finite simple games \cite{Isbell}; complete simple
games are those for which the resulting preorder is total, and their player
types are its equivalence classes \cite{CarrerasFreixas}.  Carreras and Owen
studied automorphisms in that setting \cite{CarrerasOwen}.  Klivans used a
generalized vicinal preorder to characterize finite shifted simplicial
complexes by totality \cite{Klivans}.  Kalai also proved a finite rigidity
statement for shifted families in the stronger language of weak isomorphism
\cite[Proposition~4.2]{Kalai}.

On the Banach-space side, Brech--Ferenczi--Tcaciuc supplied the decisive
signed-permutation reduction \cite{BFT}; Brech--Pi\~na later established the
rigidity of hereditary spreading families under permutation
\cite{BP}; and the 2022 survey still recorded the explicit-isometry-group
problem \cite{ABsurvey}.

A targeted search of these references, their citation trails, arXiv, and
recent work through 21 June 2026 did not locate the exact statement of
Theorems~\ref{thm:aut} and \ref{thm:main-isometry} for countably infinite
spreading families and combinatorial Banach spaces.  The finite game-theoretic
and shifted-complex literature is a substantial antecedent, so this search
result should be read only as a careful status report, not as a guarantee of
novelty.

\section{Conclusion}

The supplied threshold argument was correct, but its invariant was revealing
a more general object: the replacement preorder.  Spreading makes that
preorder total, the well-order of $\N$ prevents automorphisms from moving its
classes, and mutual replaceability makes every permutation inside a class
legal.  This yields the complete formula
\[
  \Aut(\F)=\prod_{C\in\Ccal_{\F}}\Sym(C)
\]
for every spreading family and, after the theorem of
Brech--Ferenczi--Tcaciuc, the full signed formula for every regular
combinatorial Banach space and its dual.

\begin{thebibliography}{99}

\bibitem{ABsurvey}
L. Antunes and K. Beanland,
\emph{Surjective isometries on Banach sequence spaces: A survey},
Concrete Operators \textbf{9} (2022), 19--40.
\href{https://doi.org/10.1515/conop-2022-0125}{doi:10.1515/conop-2022-0125}.

\bibitem{BFT}
C. Brech, V. Ferenczi, and A. Tcaciuc,
\emph{Isometries of combinatorial Banach spaces},
Proc. Amer. Math. Soc. \textbf{148} (2020), 4845--4854.
\href{https://doi.org/10.1090/proc/15122}{doi:10.1090/proc/15122}.

\bibitem{BP}
C. Brech and C. Pi\~na,
\emph{Banach--Stone-like results for combinatorial Banach spaces},
Ann. Pure Appl. Logic \textbf{172} (2021), 102989.
\href{https://doi.org/10.1016/j.apal.2021.102989}{doi:10.1016/j.apal.2021.102989}.

\bibitem{CarrerasFreixas}
F. Carreras and J. Freixas,
\emph{Complete simple games},
Math. Social Sci. \textbf{32} (1996), 139--155.

\bibitem{CarrerasOwen}
F. Carreras and G. Owen,
\emph{Automorphisms and weighted values},
Internat. J. Game Theory \textbf{26} (1997), 1--10.
\href{https://doi.org/10.1007/BF01262508}{doi:10.1007/BF01262508}.

\bibitem{Isbell}
J. R. Isbell,
\emph{A class of simple games},
Duke Math. J. \textbf{25} (1958), 423--439.
\href{https://doi.org/10.1215/S0012-7094-58-02537-7}{doi:10.1215/S0012-7094-58-02537-7}.

\bibitem{Kalai}
G. Kalai,
\emph{Symmetric matroids},
J. Combin. Theory Ser. B \textbf{50} (1990), 54--64.
\href{https://doi.org/10.1016/0095-8956(90)90096-I}{doi:10.1016/0095-8956(90)90096-I}.

\bibitem{Klivans}
C. J. Klivans,
\emph{Obstructions to shiftedness},
Discrete Comput. Geom. \textbf{33} (2005), 535--545.
\href{https://doi.org/10.1007/s00454-004-1103-9}{doi:10.1007/s00454-004-1103-9}.

\end{thebibliography}

\end{document}
